\documentclass[aspectratio=169]{beamer}

\usepackage[utf8]{inputenc}
\usepackage[T1]{fontenc}
\usepackage[spanish]{babel}
\usepackage{booktabs}
\usepackage{amsmath}

\usetheme{Madrid}
\usecolortheme{seahorse}

\title[Sandbox de Partículas 2D]{Sandbox de partículas interactivo 2D}
\subtitle{Comparación de ejecución en CPU secuencial vs ejecución en CUDA}
\author[Vicente Zapata González]{Vicente Zapata González\\ROL: 202073573-5}
\institute[]{
    (ILI257) Programación Paralela Aplicada\\
    Profesor: Jorge Díaz Matte
}
\date{Jueves 25 de junio de 2026}

\begin{document}

\begin{frame}
    \titlepage
\end{frame}

\begin{frame}{Motivación}
    \begin{itemize}
        \item Se implementó una simulación 2D con interfaz gráfica, donde el usuario puede observar y modificar partículas en tiempo real.
        \item Cada partícula posee posición, velocidad y material.
        \item Las partículas reaccionan al colisionar según reglas personalizadas.
        \item El objetivo principal es mostrar cómo una simulación costosa se beneficia del paralelismo en GPU.
    \end{itemize}
\end{frame}

\begin{frame}{Problema computacional}
    En cada frame, cada partícula se compara contra todas las demás:

    \[
        \text{comparaciones por frame} = N \times N
    \]

    \vspace{0.4cm}

    \begin{itemize}
        \item Con $N = 1024$: aproximadamente 1 millón de comparaciones.
        \item Con $N = 8192$: aproximadamente 67 millones de comparaciones.
        \item Con $N = 16384$: aproximadamente 268 millones de comparaciones.
    \end{itemize}

    \vspace{0.3cm}

    Aunque CPU y CUDA mantienen complejidad $O(N^2)$, la GPU puede repartir el trabajo entre miles de hilos.
\end{frame}

\begin{frame}{Arquitectura del proyecto}
    \begin{columns}
        \begin{column}{0.48\textwidth}
            \textbf{Versión CPU}
            \begin{itemize}
                \item C++ secuencial.
                \item Doble buffer de partículas.
                \item Un ciclo externo por partícula.
                \item Un ciclo interno contra todas las demás.
                \item Render con OpenGL.
            \end{itemize}
        \end{column}

        \begin{column}{0.48\textwidth}
            \textbf{Versión CUDA}
            \begin{itemize}
                \item Un hilo por partícula.
                \item Doble buffer en GPU.
                \item Kernel tiled $O(N^2)$.
                \item Uso de memoria compartida \texttt{\_\_shared\_\_}.
                \item Render con OpenGL.
            \end{itemize}
        \end{column}
    \end{columns}

    \vfill
    {\scriptsize Nota: la estructura del proyecto reutiliza como base las ideas trabajadas en los laboratorios del curso: host C++/CUDA, medición, buffers y render con OpenGL.}
\end{frame}

\begin{frame}[fragile]{CPU secuencial}
    La versión CPU calcula cada partícula una por una:

\begin{verbatim}
for (int i = 0; i < count; ++i) {
    Particle p = in[i];

    for (int j = 0; j < count; ++j) {
        if (i == j) continue;
        aplicar_interaccion_cpu(p, in[j]);
    }

    integrar_y_rebotar_cpu(p, dt);
    out[i] = p;
}
\end{verbatim}

    \begin{itemize}
        \item Es simple y directa.
        \item No usa hilos ni paralelismo explícito.
        \item Aumentar $N$ impacta muy fuerte en FPS.
    \end{itemize}
\end{frame}

\begin{frame}[fragile]{CUDA tiled}
    En CUDA, cada hilo procesa una partícula:

\begin{verbatim}
__shared__ Particle tile[Config::LOCAL_SIZE];

for (int tileStart = 0; tileStart < N; tileStart += blockDim.x) {
    tile[threadIdx.x] = inParticles[tileStart + threadIdx.x];
    __syncthreads();

    for (int k = 0; k < tileCount; ++k) {
        aplicar_interaccion_cuda(p, tile[k]);
    }

    __syncthreads();
}
\end{verbatim}

    \begin{itemize}
        \item Se reutilizan bloques de partículas en memoria compartida.
        \item Cada hilo escribe solamente su propia partícula.
        \item Se evita una condición de carrera sobre el estado global.
    \end{itemize}
\end{frame}

\begin{frame}{Interacción de materiales}
    Se definieron tres materiales con comportamiento distinto:

    \begin{itemize}
        \item \textbf{Verde}: material liviano, acelera con mayor facilidad.
        \item \textbf{Azul}: material normal, recibe una transferencia suave de velocidad.
        \item \textbf{Rojo}: material pesado, responde menos al choque.
    \end{itemize}

    \vspace{0.3cm}

    Además:

    \begin{itemize}
        \item Ninguna partícula puede quedar con velocidad cero.
        \item Se limita la velocidad máxima para mantener una simulación legible.
        \item El usuario puede seleccionar color con \texttt{1}, \texttt{2}, \texttt{3} y emitir partículas con click.
    \end{itemize}
\end{frame}

\begin{frame}{Metodología de medición}
    \begin{itemize}
        \item Tiempo por frame medido con \texttt{std::chrono}.
        \item Primeros 30 frames descartados como calentamiento.
        \item FPS promedio actualizado en consola cada 60 frames.
        \item Misma cantidad $N$ y misma lógica $O(N^2)$ en CPU y CUDA.
    \end{itemize}

    \vspace{0.4cm}

    Equipo usado:

    \begin{itemize}
        \item CPU: Intel Core i5-11400F, 6 núcleos / 12 hilos.
        \item GPU: NVIDIA GeForce RTX 4060, 3072 núcleos CUDA.
    \end{itemize}
\end{frame}

\begin{frame}{Demostración en vivo}
    \begin{center}
        \vspace{0.2cm}
        \fbox{
            \begin{minipage}[c][5.0cm][c]{0.82\textwidth}
                \centering
                Ejecución del sandbox durante la presentación\\[0.25cm]
                Comparación directa: CPU secuencial vs CUDA\\[0.25cm]
                Cantidad de partículas ajustada según el equipo disponible
            \end{minipage}
        }

        \vspace{0.25cm}
    \end{center}
\end{frame}

\begin{frame}{Resultados}
    \begin{table}
        \centering
        \small
        \begin{tabular}{rrrr}
            \toprule
            $N$ & CPU FPS & CUDA FPS & Speedup \\
            \midrule
            1024  & 383 & 600 & 1.57x \\
            2048  & 147 & 552 & 3.76x \\
            4096  & 46  & 530 & 11.64x \\
            8192  & 12  & 369 & 29.93x \\
            16384 & 3   & 338 & 110.37x \\
            \bottomrule
        \end{tabular}
    \end{table}

    {\scriptsize $N$ corresponde a la cantidad total de partículas simuladas.}

    \vspace{0.25cm}

    \begin{itemize}
        \item Al aumentar $N$, el costo cuadrático se vuelve dominante en CPU.
        \item En $N = 16384$, CPU baja a aproximadamente 3 FPS.
        \item CUDA mantiene más de 300 FPS incluso con 268 millones de comparaciones por frame.
        \item El speedup crece desde 1.57x hasta 110.37x.
    \end{itemize}

    \vspace{0.3cm}

    Esto muestra que la GPU no cambia la complejidad teórica, pero sí cambia drásticamente la capacidad práctica de procesar el trabajo.
\end{frame}

\begin{frame}{Dificultades encontradas}
    \begin{itemize}
        \item OpenCL en WSL no expuso correctamente la GPU NVIDIA.
        \item Se decidió usar CUDA como backend GPU principal.
        \item Fue necesario mantener reglas equivalentes en CPU y CUDA.
        \item Se debió evitar que varias hebras escribieran la misma partícula.
        \item Hubo que ajustar velocidades para evitar simulaciones visualmente caóticas.
    \end{itemize}
\end{frame}

\begin{frame}{Cumplimiento de la propuesta}
    La propuesta indicaba que una simulación de partículas se beneficiaría de una implementación paralela porque:

    \begin{itemize}
        \item muchas partículas deben actualizar posición, velocidad y estado;
        \item la comparación directa entre partículas escala muy rápido;
        \item muchas operaciones pueden calcularse de manera independiente;
        \item la GPU puede procesar múltiples partículas simultáneamente.
    \end{itemize}

    \vspace{0.3cm}

    Los resultados confirman esa hipótesis: al crecer $N$, la CPU secuencial cae fuertemente, mientras CUDA mantiene una simulación fluida.
\end{frame}

\begin{frame}{Conclusión}
    \begin{itemize}
        \item Se implementó un sandbox 2D interactivo con reglas personalizadas.
        \item Se comparó CPU secuencial contra CUDA usando la misma lógica $O(N^2)$.
        \item CUDA obtuvo mejoras crecientes al aumentar el número de partículas.
        \item La mayor mejora medida fue de aproximadamente 110x con $N = 16384$.
        \item El proyecto cumple el objetivo de demostrar visual y cuantitativamente el beneficio del paralelismo en GPU.
    \end{itemize}
\end{frame}

\begin{frame}
    \centering
    \vspace{1.0cm}
    {\Huge Gracias por su atención}
\end{frame}

\end{document}
